\scanpage{36}
est cependant possible que l'injectivité de (*) reste valable sans hyp. de caractéristique -- à examiner de plus près...

Revenant aux conditions du cor. 1, b), en car. $p>0$, et travaillant avec les groupes fondamentaux premiers à $p$, on trouve que le $u$ dans 1°) est un point \inserted{de $G(K)$} d'ordre une puissance de $p$. Si $G$ est un tore (il suffit pour \unsure{cela|celui} que les ``fibres'' dans la variété élémentaire $Y$ soient \unsure{dé|de} genre 0) on en conclut encore que $u=0$, i.e. $f=g$. Dans le cas général, la démonstration arithmétique prouve \inserted{(sauf erreur)} ceci : si $x\in X$, \deleted{et} \inserted{d'où} $f(x)\inserted{=P}$, $g(x)=Q$ dans $Y$, non seulement $Q-P$ est un élément de $p$-torsion de la \unsure{jacobienne|Jacobienne} gén. $G_Y$ de $Y$, pour tt rev. \inserted{étale} galoisien fini $Y'$ de $X$, \inserted{de degré premier à $p$\unsure{,|,,}} pointé par un $P'$ au dessus de $P$, il existe $Q'$ au dessus de $Q$ tel que $Q'-P'$ soit un élément de ? $p$-torsion de $G_{Y'}$.

\deleted{en \unsure{prenant|passant} comme \unsure{revêtements|relèvements} les \inserted{rev. abéliens, $=$ \unsure{induits} par les} isogénies $\unsure{n_G|{}_nG|nG}$ ($n$ premier à \unsure{$G$|$p$}), il existe un $Q_n\in G$ tel que $\unsure{{}_n n\cdot Q_n|n\cdot Q_n}=Q$, et que $Q_n$}
\unsure{Il|Je} se pourrait qu'on puisse en déduire $P=Q$

\margin{Ceci montre que le \unsure{th.|thm.} 2 reste vrai, sur un corps de base de car. $p>0$, en travaillant avec les groupes fondamentaux premiers à $p$ -- car là dans la réduction à un énoncé géométrique, on voit $X$, de sorte qu'on peut supposer que $Y$ est un \uline{produit} de courbes de genre zéro. Dans ce cas d'ailleurs, on est ramené au cas où $Y$ est lui $=$ une telle courbe,...}
\margin{35}
